\section{Introduction}


Skilled immigration sits at the center of several debates in labor economics. One set of questions concerns the effects of high-skill immigration on firms, innovation, and the employment of native workers. Another concerns the private consequences of immigration policy for the migrants whose careers are directly shaped by visa rules. The H-1B program is especially important in both debates because it is the main capped channel through which U.S. employers hire foreign workers in specialty occupations. Yet while a large literature studies how changes in H-1B access affect firms and local labor markets, much less is known about what happens to the individual worker who wins or loses the H-1B lottery. That missing worker-level margin matters for both policy and welfare. If losing the lottery typically forces workers to leave the sponsoring employer or to leave the United States altogether, then the lottery creates large private disruptions even if firm-level responses are muted. If instead many losers remain through alternative channels, then the H-1B cap may operate less as a sharp entry barrier and more as a source of delay and uncertainty.

Estimating those individual effects is difficult for two reasons. First, visa receipt is usually endogenous to labor-market outcomes. Firms apply for workers they value, and workers who obtain visas differ systematically from those who do not. The H-1B cap lottery offers a natural solution to that problem, but only if one can observe both winners and losers. Second, the administrative data that record H-1B applications are typically deidentified. In the absence of names or stable person identifiers, it is challenging to follow individual applicants into longitudinal labor-market data, especially for lottery losers, who may disappear from standard U.S. administrative sources if they leave the country or never transition onto H-1B status. Existing work has therefore focused largely on firm-level outcomes or on selected employer panels rather than on the individual applicant.

This paper overcomes both obstacles. I combine the universe of recent cap-subject H-1B lottery records obtained through FOIA requests with LinkedIn-based individual work-history data from Revelio Labs. The empirical leverage comes from a change in the institutional environment surrounding the H-1B program. In the older view of the program, many applicants were abroad at the time of filing, so employers only observed workers inside the firm if the petition succeeded. In the modern setting, a large share of applicants are already working in the United States before the lottery, most commonly under F-1 Optional Practical Training (OPT). For this group, the sponsoring firm often employs the worker before it enters the lottery. That fact makes it possible to condition on pre-lottery employment, define a credible candidate pool inside the sponsoring firm, and link deidentified H-1B applicants to individual LinkedIn profiles using probabilistic record linkage.

The first contribution of the paper is methodological. I develop a linkage procedure tailored to deidentified immigration records. The procedure begins with an employer crosswalk between the administrative H-1B files and Revelio firms, combining exact and fuzzy company-name linkage with LLM-assisted review of ambiguous cases. On the Revelio side, I then impute the sparse fields needed for individual linkage---country of origin, year of birth, sex, and advanced-degree status---using education histories, firm histories, and name-based models. The resulting applicant-candidate scores are converted into within-applicant match weights, which I carry forward into the outcome analysis instead of forcing a single deterministic match. This feature is useful in precisely the setting studied here, where no single identifier is decisive and where researchers would like to preserve uncertainty rather than hide it.

The second contribution is substantive. To the best of my knowledge, this is the first paper to estimate the causal effect of winning the H-1B lottery on the individual worker using application-level administrative data on both winners and losers. My empirical design compares applicants who were selected and not selected in the lottery, conditioning on advanced-degree-exempt status because U.S. master's and doctoral degree holders face a different lottery path. In the linked sample, I collapse candidate-level outcomes to the applicant level using the probabilistic match weights and estimate specifications with and without firm-by-lottery-year fixed effects. Identification relies on the standard lottery assumption that, conditional on advanced-degree status, selection is independent of applicant potential outcomes.

The current estimates point to a consistent but modest effect of winning. Lottery winners are more likely to remain with the sponsoring employer and more likely to be observed working in the United States in the years after the lottery. In the current draft, the effect on same-firm retention two years after the lottery is about 0.9 percentage points, and the effect on remaining in the United States is about 2.8 percentage points. These effects are economically meaningful but small relative to baseline attachment. The magnitudes appear to reflect two forces. First, probabilistic linkage inevitably introduces measurement error, which attenuates treatment effects toward zero. Second, the small estimates seem to be substantive rather than purely mechanical: even in high-quality matched subsamples, many lottery losers remain at the firm or in the United States through other channels, including STEM OPT extensions, subsequent lottery attempts, cap-gap protection, or alternative visa paths. The effects are correspondingly larger for workers who appear to be closer to the end of their OPT runway.

These findings speak to several literatures. A long tradition studies the aggregate and firm-level consequences of skilled immigration policy, including innovation, productivity, and employment effects of changes in H-1B access \citep{boundetal2017, kerrlincoln2010, doranetal2022, maydaetal2020, mahajanetal2024}. A related line of work emphasizes how immigration rules can shape firms' outside options and the global geography of work, for example by inducing offshoring or affecting destination choices outside the United States \citep{glennon2024, khannamorales2017, brinattiguo2023}. Another literature highlights that temporary work visas may generate employer-specific rents by restricting worker mobility and tying legal status to the sponsoring employer \citep{kimpei2022}. This paper complements those studies by shifting the unit of analysis from the firm to the worker. It asks not how an additional H-1B affects the sponsoring firm, but how winning the lottery changes the path of the individual applicant.

The results also suggest a different interpretation of the modern H-1B program for a large and growing subset of applicants. Among U.S.-trained workers already employed under OPT, the H-1B often functions less as a mechanism for initial entry and more as a mechanism for retention and longer-horizon attachment. That distinction matters for how one interprets both efficiency and distributional concerns. If many applicants are already present and productive in the United States before the lottery, then the main effect of lottery loss may be to weaken firm attachment, reduce employment stability, or accelerate return migration only at the margin. At the same time, the modest average effects should not be read as evidence that the lottery is unimportant. A policy that creates substantial uncertainty, repeated reapplication costs, and employer dependence can have large welfare consequences even when average reduced-form employment responses are numerically small.

The remainder of the paper proceeds as follows. Section \ref{sec:context} describes the institutional setting, with particular attention to the H-1B lottery and the role of OPT. Section \ref{sec:data} introduces the administrative H-1B records and the Revelio outcome data. Section \ref{sec:linkage} describes the probabilistic linkage procedure and the construction of the linked applicant panel. Section \ref{sec:results} lays out the empirical design and organizes the results. Section \ref{sec:conclusion} concludes.
